\documentclass{article}
\usepackage[UTF8]{ctex}
\usepackage{amsmath}
\usepackage{amssymb}
\usepackage{graphicx}
\usepackage{geometry}
\geometry{a4paper, margin=2cm}

\newcommand{\problemname}[1]{\begin{center}\Large\bfseries #1\end{center}}
\usepackage{enumitem}

\begin{document}
\problemname{肥胖问题}

\subsection*{题目描述}
用公式 x = 100000 × m / (h × h) 计算胖瘦指数 x（结果只取整数部分，不考虑四舍五入），其中 m 为体重（公斤），h 为身高（厘米）。根据 x 进行体型判定并输出对应结果：x 小于 180 输出 Underweight；x 大于等于 180 且小于 240 输出 Normal；x 大于等于 240 输出 Overweight。

\subsection*{输入格式}
一行两个整数 m 和 h，用空格分隔，满足 40 ≤ m ≤ 120，140 ≤ h ≤ 200。

\subsection*{输出格式}
一行字符串：Underweight、Normal 或 Overweight（每次输出后换行）。

\subsection*{样例输入}
\begin{verbatim}
45 160
\end{verbatim}

\subsection*{样例输出}
\begin{verbatim}
Underweight
\end{verbatim}


\end{document}
